\documentclass[11pt,a4paper]{article}

% ------------------------------------------------------------------ %
% Packages
% ------------------------------------------------------------------ %
\usepackage[margin=2.2cm]{geometry}
\usepackage{amsmath,amssymb,mathtools}
\usepackage{booktabs}
\usepackage{longtable}
\usepackage{array}
\usepackage{xcolor}
\usepackage{listings}
\usepackage{siunitx}
\usepackage{microtype}
\usepackage{enumitem}
\usepackage{algorithm}
\usepackage{algpseudocode}
\usepackage[hidelinks]{hyperref}

\algrenewcommand\algorithmicrequire{\textbf{Input:}}
\algrenewcommand\algorithmicensure{\textbf{Output:}}
\newcommand{\pyconst}[1]{\textcolor{black!60!blue}{\texttt{\footnotesize #1}}}
\newcommand{\dictkey}[1]{\textcolor{black!60!green}{\texttt{\footnotesize #1}}}

% ------------------------------------------------------------------ %
% Listing style for parameter snippets
% ------------------------------------------------------------------ %
\lstdefinestyle{cfg}{
  basicstyle=\ttfamily\footnotesize,
  breaklines=true,
  columns=fullflexible,
  frame=single,
  rulecolor=\color{black!30},
  showstringspaces=false,
}
\lstset{style=cfg}

% ------------------------------------------------------------------ %
% Convenience macros
% ------------------------------------------------------------------ %
% \Vm and \Iion are spelt without any built-in subscript so they can be
% safely subscripted at call sites (e.g. \Vmsub{c}). This avoids the
% "double subscript" pitfall when one writes \Vm_{c}.
\newcommand{\Vm}{V_{m}}
\newcommand{\Vmsub}[1]{V_{m,#1}}
\newcommand{\Iion}{I_{\text{ion}}}
\newcommand{\Iionsub}[1]{I_{\text{ion},#1}}
\newcommand{\Iext}{I_{\text{ext}}}
\newcommand{\Dten}{\mathbf{D}}
\newcommand{\grad}{\nabla}
\newcommand{\divg}{\nabla\!\cdot}
\newcommand{\Aimp}{\mathsf{A}}
\newcommand{\Bimp}{\mathsf{B}}
\newcommand{\Mmass}{\mathsf{M}}
\newcommand{\Kstif}{\mathsf{K}}
\newcommand{\Lop}{\mathsf{L}}
\newcommand{\code}[1]{\texttt{\small #1}}

\title{Solver Specification:\\
       \texorpdfstring{\texttt{highOrderElectroActivationFoamImplicitJfNK\_PETSc}}{highOrderElectroActivationFoamImplicitJfNK\_PETSc}\\
       \large A high-order, implicit, Jacobian-free Newton--Krylov
       monodomain solver in OpenFOAM with PETSc back-end}
\author{Pablo}
\date{\today}

\begin{document}
\maketitle

\begin{abstract}
This report specifies the algorithms implemented in the OpenFOAM solver
\code{highOrderElectroActivationFoamImplicitJfNK\_PETSc}. The solver
advances the monodomain equations of cardiac electrophysiology coupled
to the ten Tusscher--Noble--Noble--Panfilov (TNNP, 2004) ionic cell
model. The diffusion operator is discretised with the cell-centred
finite-volume framework of OpenFOAM, optionally upgraded to a Local
Radial-Extrapolation (LRE) high-order reconstruction; the ionic source
is evaluated either at cell centres or at LRE cell quadrature points;
time integration follows a $\theta$-scheme (backward Euler or
Crank--Nicolson); the resulting nonlinear coupled system is solved by a
Jacobian-Free Newton--Krylov (JFNK) method with Armijo backtracking line
search, linear extrapolation initial guess and a physical-bracket clamp
on the input to the ionic ODE; the inner linear system is solved by
PETSc's KSP framework (GMRES with modified Gram--Schmidt by default);
the per-cell ionic ODE uses OpenFOAM's adaptive RKF45 integrator with
configurable tolerances; activation times are tracked by linear
interpolation across the first threshold crossing and benchmark points
are sampled by inverse-distance weighting. All algorithmic parameters
are exposed through the \code{constant/spatialIntegrationProperties}
and \code{constant/timeIntegrationProperties} dictionaries, and the
companion Python driver \code{run\_convergence.py} generates these
dictionaries from a single configuration file.
\end{abstract}

\tableofcontents
\bigskip

% ================================================================== %
\section{Notation}
% ================================================================== %

\begin{center}
\begin{tabular}{@{}lll@{}}
\toprule
Symbol & Meaning & Units \\
\midrule
$\Omega \subset \mathbb{R}^{d}$ & Spatial domain ($d \in \{2,3\}$) & \si{\meter} \\
$\partial\Omega$ & Insulated boundary (zero-flux Neumann) & --- \\
$t \in (0, T]$ & Time & \si{\second} \\
$\Vm(\mathbf{x}, t)$ & Transmembrane voltage & \si{\volt} \\
$\mathbf{s}(\mathbf{x}, t) \in \mathbb{R}^{17}$ & TNNP ionic states & mixed \\
$\Iion(\Vm, \mathbf{s})$ & Total ionic current density & \si{\volt\per\second} \\
$\Iext(\mathbf{x}, t)$ & External stimulus current density & \si{\ampere\per\meter\cubed} \\
$\chi$ & Membrane surface-to-volume ratio & \si{\per\meter} \\
$C_{m}$ & Membrane capacitance per unit area & \si{\farad\per\meter\squared} \\
$\Dten(\mathbf{x})$ & Bulk conductivity tensor & \si{\siemens\per\meter} \\
$\theta \in \{0.5, 1\}$ & Time-stepping parameter (CN or BE) & --- \\
$\Delta t$ & Outer time-step length & \si{\second} \\
$N_{c}$ & Number of mesh cells & --- \\
$\mathsf{V}^{n} \in \mathbb{R}^{N_{c}}$ & Vector of cell-averaged $\Vm$ at time $t^{n}$ & \si{\volt} \\
$\Mmass, \Kstif, \Lop$ & Mass, stiffness, scaled diffusion matrix & --- \\
$\Aimp, \Bimp$ & Implicit LHS / RHS operators of the $\theta$-scheme & --- \\
$\mathsf{F}(\mathsf{x})$ & Discrete nonlinear residual to be zeroed & --- \\
\bottomrule
\end{tabular}
\end{center}

% ================================================================== %
\section{Governing equations}
% ================================================================== %

\subsection{Monodomain PDE}

The bulk voltage is propagated by the standard monodomain reaction--diffusion
equation \cite{ttp2004,niederer2012}:
\begin{equation}
  \chi \, C_{m} \, \frac{\partial \Vm}{\partial t}
   \;=\;
   \divg\!\bigl( \Dten \grad \Vm \bigr)
   - \Iion(\Vm, \mathbf{s})
   + \Iext
   \quad \text{in } \Omega \times (0, T],
   \label{eq:monodomain}
\end{equation}
subject to insulated boundaries
\begin{equation}
  \mathbf{n} \cdot \bigl(\Dten \grad \Vm\bigr) \;=\; 0
   \quad \text{on } \partial\Omega \times (0, T],
\end{equation}
and initial condition $\Vm(\mathbf{x}, 0) = -\SI{84}{\milli\volt}$.
The conductivity tensor $\Dten$ is assumed symmetric positive-definite
and either uniform or read as a \code{volTensorField} from disk.

Dividing~\eqref{eq:monodomain} by $\chi C_{m}$ gives the form actually
used inside the solver,
\begin{equation}
  \frac{\partial \Vm}{\partial t}
   \;=\;
   \mathcal{L} \Vm
   \;+\;
   s(\Vm, \mathbf{s}),
   \qquad
   \mathcal{L} \;=\; \frac{1}{\chi C_{m}}\,\divg(\Dten \grad \cdot),
   \quad
   s \;=\; -\frac{\Iion}{\chi C_{m}} + \frac{\Iext}{\chi C_{m}}.
   \label{eq:semi-discrete}
\end{equation}

\subsection{Ionic ODE}

At every spatial sampling point (cell centre or LRE cell quadrature
point) the ionic state vector $\mathbf{s} \in \mathbb{R}^{17}$ evolves
according to the TNNP~2004 system
\begin{equation}
  \frac{\mathrm{d}\mathbf{s}}{\mathrm{d}t}
   \;=\;
   \mathbf{f}_{\text{TNNP}}\!\bigl(\Vm, \mathbf{s}; \texttt{tissue}\bigr),
   \qquad \mathbf{s}(0) = \mathbf{s}_{0},
   \label{eq:tnnp}
\end{equation}
where $\Vm$ is supplied externally by the PDE solver, and
\code{tissue} selects one of \code{epicardialCells},
\code{mCells} or \code{endocardialCells}. The 17 states tracked in
\code{TNNP\_2004Names.H} are
\[
  \mathbf{s} = (V,\ K_{i},\ Na_{i},\ Ca_{i},\ Xr_{1},\ Xr_{2},\ Xs,\ m,\ h,\ j,\ d,\ f,\ fCa,\ s,\ r,\ g,\ Ca_{SR}).
\]
The total ionic current $\Iion$ returned to the PDE is the algebraic
quantity \code{Iion\_cm} (sum of $i_{Na}, i_{K1}, i_{to}, i_{Kr},
i_{Ks}, i_{CaL}, i_{NaK}, i_{NaCa}, i_{pK}, i_{pCa}, i_{bNa},
i_{bCa}$). See the original publication \cite{ttp2004} for the
algebraic definitions of each current.

\subsection{External stimulus}

The external current $\Iext$ is supplied as a time-windowed,
region-windowed step,
\begin{equation}
  \Iext(\mathbf{x}, t) \;=\;
  \begin{cases}
    I_{0} & \mathbf{x} \in \Omega_{\text{stim}}^{(b)},
            \;\; t \in [t_{0}^{(b)},\, t_{0}^{(b)} + \tau],
            \;\; b = 1, \dots, B, \\
    0     & \text{otherwise},
  \end{cases}
\end{equation}
with stimulus boxes $\Omega_{\text{stim}}^{(b)}$, start times
$t_{0}^{(b)}$, common amplitude $I_{0}$ and duration $\tau$ read from
\code{constant/stimulusProtocol}. Default Niederer-benchmark values are
$I_{0} = \SI{50000}{\ampere\per\meter\cubed}$ and
$\tau = \SI{2}{\milli\second}$.

% ================================================================== %
\section{Time discretisation: the $\theta$-scheme}
% ================================================================== %

The semi-discrete equation~\eqref{eq:semi-discrete} is advanced from
$t^{n}$ to $t^{n+1} = t^{n} + \Delta t$ by a one-parameter family
\begin{equation}
  \frac{\mathsf{V}^{n+1} - \mathsf{V}^{n}}{\Delta t}
   \;=\;
   \theta\, \Lop\, \mathsf{V}^{n+1}
   + (1-\theta)\, \Lop\, \mathsf{V}^{n}
   + \theta\, \mathbf{s}^{n+1}
   + (1-\theta)\, \mathbf{s}^{n}.
   \label{eq:theta}
\end{equation}
Two values of $\theta$ are selectable by the dictionary key
\code{implicitScheme}:
\begin{center}
\begin{tabular}{@{}lcl@{}}
\toprule
\code{implicitScheme} & $\theta$ & Scheme \\
\midrule
\code{backwardEuler} & $1$   & First-order, L-stable, fully implicit \\
\code{crankNicolson} & $1/2$ & Second-order, A-stable (not L-stable) \\
\bottomrule
\end{tabular}
\end{center}

Multiplying~\eqref{eq:theta} by $\Delta t^{-1}$ on both sides yields the
linear system actually assembled,
\begin{align}
  \bigl[\Mmass/\Delta t \;-\; \theta\, \Lop\bigr]\, \mathsf{V}^{n+1}
  &\;=\;
  \bigl[\Mmass/\Delta t \;+\; (1-\theta)\, \Lop\bigr]\, \mathsf{V}^{n} \notag\\
  &\quad+\; \theta\, \mathbf{s}^{n+1}
        +\; (1-\theta)\, \mathbf{s}^{n}.
   \label{eq:implicit-system}
\end{align}
We name the two PETSc \code{Mat} operators in~\eqref{eq:implicit-system}
\begin{equation}
  \Aimp \;=\; \Mmass/\Delta t \,-\, \theta\, \Lop,
   \qquad
  \Bimp \;=\; \Mmass/\Delta t \,+\, (1-\theta)\, \Lop.
   \label{eq:AB}
\end{equation}
They are rebuilt at every time step by duplicating $\Mmass$, scaling by
$\Delta t^{-1}$ and adding $\pm \theta \Lop$ via \code{MatAXPY} with the
\code{DIFFERENT\_NONZERO\_PATTERN} flag (since $\Mmass$ and $\Lop$ have
disjoint sparsity patterns in general). When $\theta = 1$ (BE) the
explicit half of the diffusion vanishes and the second \code{MatAXPY}
call is skipped, leaving $\Bimp = \Mmass/\Delta t$.

The scaled diffusion operator $\Lop$ is itself derived from the
stiffness matrix $\Kstif$ by
\begin{equation}
  \Lop \;=\; \frac{1}{\chi\, C_{m}}\, \Kstif,
\end{equation}
assembled once at solver start-up by \code{MatDuplicate} +
\code{MatScale} on a deep copy of $\Kstif$.

% ================================================================== %
\section{Spatial discretisation}
% ================================================================== %

The mesh is the standard OpenFOAM \code{fvMesh}. All operators below
return PETSc \code{Mat} objects of type \code{MATSEQAIJ} with
preallocated non-zero counts per row.

\subsection{Standard orthogonal FV Laplacian}

\code{assembleStandardOrthogonalStiffnessMatrix} computes the
classical 5/7-point Laplacian. For each internal face $f$ between cells
$\text{own}$ and $\text{nei}$ with outward (from own) face-area vector
$\mathbf{S}_{f}$ and cell-centre separation $\mathbf{d}_{f} =
\mathbf{C}_{\text{nei}} - \mathbf{C}_{\text{own}}$:
\begin{equation}
  a_{f} \;=\;
   \frac{ \lvert \mathbf{S}_{f} \rvert\,
          \bigl( \mathbf{n}_{f} \cdot \Dten_{f} \cdot \hat{\mathbf{d}}_{f} \bigr) }
        { \lvert \mathbf{d}_{f} \rvert },
   \quad
   \Dten_{f} = \tfrac{1}{2}\,(\Dten_{\text{own}} + \Dten_{\text{nei}}),
   \quad
   \hat{\mathbf{d}}_{f} = \mathbf{d}_{f}/\lvert \mathbf{d}_{f} \rvert.
\end{equation}
Per face the four matrix contributions assembled with \code{ADD\_VALUES}
into $\Kstif$ are
\begin{equation}
\begin{aligned}
  \Kstif[\text{own},\text{own}] &\;\mathrel{+}=\; -a_{f}/V_{\text{own}}, &
  \Kstif[\text{own},\text{nei}] &\;\mathrel{+}=\; +a_{f}/V_{\text{own}}, \\
  \Kstif[\text{nei},\text{own}] &\;\mathrel{+}=\; +a_{f}/V_{\text{nei}}, &
  \Kstif[\text{nei},\text{nei}] &\;\mathrel{+}=\; -a_{f}/V_{\text{nei}}.
\end{aligned}
\end{equation}
Faces on a \code{zeroGradient} or \code{empty} patch contribute zero
flux. Faces on a \code{fixedValue} or \code{fixedVoltage} patch
contribute only to the owner diagonal; the prescribed value enters the
right-hand side via OpenFOAM's boundary-correction mechanism. The
sparsity is preallocated at 7 non-zeros per row, sufficient for any
hexahedral / unstructured mesh up to six face neighbours.

\subsection{High-order LRE diffusion}

When \code{useHighOrder\_Vm} is \code{true} and the geometric dimension
exceeds one, \code{assembleHighOrderStiffnessMatrix} is used instead.
It performs Gauss quadrature on each face with reconstructed gradients
provided by the LRE library:
\begin{equation}
  \Phi_{f} \;=\;
   \sum_{q \in \mathcal{Q}_{f}}
       w_{q}\, \lvert \mathbf{S}_{f} \rvert\,
       \mathbf{n}_{f} \cdot \Dten_{f} \cdot
       \underbrace{\Bigl(
         \sum_{j \in \mathcal{S}_{f}}
            \mathbf{g}_{f,q,j}\, \mathsf{V}_{j}
       \Bigr)}_{\grad \Vm \big|_{q}},
   \label{eq:HO-flux}
\end{equation}
where $\mathcal{Q}_{f}$ are the face quadrature points,
$\mathcal{S}_{f}$ is the LRE face stencil (typically $15$--$50$ cells)
and $\mathbf{g}_{f,q,j}$ are the \code{QRGradFaceGPCoeffs} returned by
the LRE object. The contribution of each $\mathsf{V}_{j}$ to the flux
is added to row \code{own} with weight $+\Phi_{f,j}/V_{\text{own}}$
and to row \code{nei} with weight $-\Phi_{f,j}/V_{\text{nei}}$, giving
the expected negative semi-definite Laplacian. Preallocation is
60~non-zeros per row.

\subsection{Mass matrices}

Two mass matrices are available:

\paragraph{Lumped (default).} \code{assembleDiagonalMassMatrix} writes
$\Mmass = I$ on the cell-centre basis,
\begin{equation}
  \Mmass[i, i] \;=\; 1, \qquad i = 1, \dots, N_{c},
\end{equation}
preallocated at one non-zero per row.

\paragraph{Consistent high-order.}
\code{assembleConsistentMassMatrixHO} fills a multi-stencil mass
matrix by Gauss quadrature on each cell. For every cell $i$ and every
stencil neighbour $j \in \mathcal{S}_{i} \cup \{i\}$,
\begin{equation}
  \Mmass[i, j] \;=\;
   \sum_{q \in \mathcal{Q}_{i}}
       \frac{w_{q}}{\sum_{q'} w_{q'}}\,
       \phi_{j}(\mathbf{x}_{q}),
   \label{eq:HO-mass}
\end{equation}
where $\phi_{j}$ is the LRE Taylor reconstruction basis function of
neighbour $j$ at quadrature point $\mathbf{x}_{q}$:
\begin{equation}
  \phi_{j}(\mathbf{x}_{q})
   \;=\;
   \delta_{ij}
   \;+\; \mathbf{g}_{i,j} \cdot \mathbf{d}_{q}
   \;+\; \tfrac{1}{2}\, \mathbf{H}_{i,j} \!:\! (\mathbf{d}_{q} \otimes \mathbf{d}_{q})
   \;+\; \tfrac{1}{6}\, \mathbf{T}^{(3)}_{i,j} \!\vdots\!
       (\mathbf{d}_{q} \otimes \mathbf{d}_{q} \otimes \mathbf{d}_{q}),
   \label{eq:taylor}
\end{equation}
with $\mathbf{d}_{q} = \mathbf{x}_{q} - \mathbf{C}_{i}$. The gradient,
Hessian and 3rd-order tensor coefficients are
\code{QRGradCoeffs}, \code{cellHessianCoeffs} and
\code{cellThirdDerivCoeffs} respectively, and the cubic term is
contracted by \code{LRE::cubicForm}. The weight normalisation in
\eqref{eq:HO-mass} guarantees that $(\Mmass\,\mathbf{1})_{i} = 1$ for
constant fields (consistency on constants is exact at any order).

A safety guard at line 1500 of the source forces $\Mmass$ to the lumped
form whenever \code{useHighOrder\_Vm = false} or $d \le 1$, regardless
of \code{massMatrix}, with an informational message.

\subsection{High-order reconstruction at ionic quadrature points}

When \code{useHighOrder\_Iion} is \code{true} the ionic ODE is
integrated at cell quadrature points rather than at cell centres. The
voltage is reconstructed there by the same Taylor expansion as
in~\eqref{eq:taylor},
\begin{equation}
  \Vm(\mathbf{x}_{q})
   \;\approx\;
   \Vmsub{c} + (\grad \Vm)_{c} \cdot \mathbf{d}_{q}
   + \tfrac{1}{2}\, \mathbf{H}_{c} \!:\! (\mathbf{d}_{q} \otimes \mathbf{d}_{q})
   + \tfrac{1}{6}\, \mathbf{T}^{(3)}_{c} \!\vdots\!
       (\mathbf{d}_{q} \otimes \mathbf{d}_{q} \otimes \mathbf{d}_{q}),
   \label{eq:HO-vm}
\end{equation}
implemented in \code{reconstructVmAtIionIntegrationPoints}. The order
of \eqref{eq:HO-vm} is the one of the \emph{Vm} LRE object,
independently of the LRE object used for $\Iion$ integration. If
\code{useHighOrder\_Vm} is \code{false} the linear fallback
$\Vmsub{c} + (\grad\Vm)_{c}\!\cdot\! \mathbf{d}_{q}$ is used with
$(\grad\Vm)_{c}$ from the standard OpenFOAM \code{fvc::grad}.

\subsection{High-order cell averaging of $\Iion$}

After the ODE solve, \code{averageIionIntegrationPointsToCells}
reduces the per-quadrature-point ionic currents to cell averages by
\begin{equation}
  \Iionsub{i} \;=\; \frac{\sum_{q \in \mathcal{Q}_{i}} w_{q}\, \Iionsub{q}}
                       {\sum_{q \in \mathcal{Q}_{i}} w_{q}},
\end{equation}
which is the natural counterpart of the Gauss reconstruction
in~\eqref{eq:HO-vm}.

% ================================================================== %
\section{Nonlinear solver}
% ================================================================== %

The fully discrete equation~\eqref{eq:implicit-system} is nonlinear in
$\mathsf{V}^{n+1}$ through $\Iion(\mathsf{V}^{n+1}, \mathbf{s}^{n+1})$,
where $\mathbf{s}^{n+1}$ is itself obtained by integrating
\eqref{eq:tnnp} from $t^{n}$ to $t^{n+1}$ with input
$\mathsf{V}^{n+1}$. We define the discrete residual
\begin{equation}
  \mathsf{F}(\mathsf{x})
   \;=\;
   \Aimp\, \mathsf{x}
   \;-\;
   \mathsf{r}\!\bigl( \Iion(\mathsf{x}) \bigr),
   \qquad
   \mathsf{x} = \mathsf{V}^{n+1},
   \label{eq:Fdef}
\end{equation}
with $\mathsf{r}(\Iion) = \Bimp\, \mathsf{V}^{n} +
\theta\, \mathbf{s}^{n+1} + (1-\theta)\, \mathbf{s}^{n}$, and seek the
unique root $\mathsf{F}(\mathsf{x}) = 0$.

\subsection{Jacobian-Free Newton--Krylov (JFNK)}

Newton's iteration is
\begin{equation}
  \mathsf{J}(\mathsf{x}_{k})\, \boldsymbol{\delta}_{k}
   \;=\; -\mathsf{F}(\mathsf{x}_{k}),
   \qquad
  \mathsf{x}_{k+1} \;=\; \mathsf{x}_{k} + \alpha_{k}\, \boldsymbol{\delta}_{k},
   \label{eq:newton}
\end{equation}
with $\mathsf{J} = \partial \mathsf{F}/\partial \mathsf{x}$ and step
length $\alpha_{k} \in (0, 1]$. The Jacobian $\mathsf{J}$ is never
assembled. Instead, every matrix--vector product
$\mathsf{J}(\mathsf{x}_{k})\, \mathbf{v}$ required by GMRES is
approximated by a forward finite difference \cite{knoll2004},
\begin{equation}
  \mathsf{J}(\mathsf{x}_{k})\, \mathbf{v}
   \;\approx\;
   \frac{\mathsf{F}(\mathsf{x}_{k} + \varepsilon\, \mathbf{v})
         - \mathsf{F}(\mathsf{x}_{k})}{\varepsilon},
   \qquad
   \varepsilon \;=\; \varepsilon_{0}\,
                     \frac{1 + \lVert \mathsf{x}_{k} \rVert_{2}}
                          {\lVert \mathbf{v} \rVert_{2}},
   \label{eq:fdjac}
\end{equation}
where $\varepsilon_{0} = $ \code{jfnkEpsilon}. The textbook choice
$\varepsilon_{0} \sim \sqrt{\text{eps}_{\text{machine}}} \approx
10^{-8}$ is optimal only when $\mathsf{F}$ is computed in exact
arithmetic; when $\mathsf{F}$ contains numerical noise of relative
magnitude $\eta$ (here $\eta \sim$ \code{relTol} of the inner RKF45
integrator), the optimum balancing truncation and noise is
$\varepsilon_{0} \approx \sqrt{\eta}$ \cite{knoll2004}. With the
default $\eta = 10^{-6}$ we use $\varepsilon_{0} = 10^{-5}$.

The wrapper \code{matVec} of equation~\eqref{eq:fdjac} is exposed to
PETSc as a \code{MatShell} so that any KSP type can act on it. The
inner GMRES solve uses modified Gram--Schmidt orthogonalisation, a
restart cycle of \code{jfnkMaxKrylovIterations} and up to
\code{jfnkMaxRestarts} restarts.

\subsection{Armijo backtracking line search}

The unconditional update $\mathsf{x}_{k+1} = \mathsf{x}_{k} +
\alpha_{0}\, \boldsymbol{\delta}_{k}$ (with $\alpha_{0} =
$~\code{nonlinearRelaxation}) is fragile near the steep $i_{Na}$
sigmoid of the cardiac AP upstroke, where the linearisation
\eqref{eq:fdjac} at $\mathsf{x}_{k}$ underestimates the true
nonlinearity. We globalise Newton by Armijo backtracking
\cite{noceWright}:
\begin{equation}
  \alpha_{k} \;=\; \alpha_{0} \cdot 2^{-\ell_{k}},
   \qquad
   \ell_{k} \;=\; \min\Bigl\{\, \ell \ge 0 \;\big|\;
       \lVert \mathsf{F}(\mathsf{x}_{k} + 2^{-\ell}\alpha_{0}
                          \boldsymbol{\delta}_{k}) \rVert_{2}^{2}
       \le \bigl(1 - 2c\, 2^{-\ell}\alpha_{0}\bigr)\,
           \lVert \mathsf{F}(\mathsf{x}_{k}) \rVert_{2}^{2}
       \Bigr\}.
   \label{eq:armijo}
\end{equation}
The sufficient-decrease constant is $c = $~\code{jfnkArmijoC}
(default $10^{-4}$). Backtracking stops when the test succeeds, after
\code{jfnkLineSearchMaxIter} trials, or when
$2^{-\ell}\alpha_{0} < $~\code{jfnkLineSearchAlphaMin}; in the latter
two cases the current trial is accepted and outer convergence
checking handles the next step. Disable with
\code{jfnkLineSearch~false} to recover the legacy fixed-damping update.

\subsection{Linear extrapolation initial guess}

Initial guess \code{jfnkInitGuessOrder} selects between
\begin{align}
  \mathsf{x}_{0} &\;=\; \mathsf{V}^{n}
       & \text{(\code{jfnkInitGuessOrder = 0})}, \\
  \mathsf{x}_{0} &\;=\; \mathsf{V}^{n} + \frac{\Delta t}{\Delta t_{\text{prev}}}
        \bigl( \mathsf{V}^{n} - \mathsf{V}^{n-1} \bigr)
       & \text{(\code{jfnkInitGuessOrder = 1}, default)}.
   \label{eq:extrap}
\end{align}
The order-1 extrapolation is cell-wise clamped to
$[\,$\code{jfnkInitGuessVmMin}, \code{jfnkInitGuessVmMax}$\,]$ to
protect Newton from starting in an unphysical bracket. On the first
time step there is no $\mathsf{V}^{n-1}$ and the solver falls back to
order 0.

\subsection{Protective Vm clamp on ODE input}

When \code{jfnkClampODEInput} is \code{true} (default), the Vm vector
passed to \code{TNNP::solveODE} inside every call to
\code{evaluateNonlinearFields} is clamped element-wise to the same
physical bracket. The clamp is applied to the working buffer
(\code{VmIntegrationPoints} for the HO-Iion path, a local
\code{scalarField} copy otherwise) so that neither \code{VmCandidate}
nor any persistent field is mutated. The kink that this introduces in
$\mathsf{F}$ outside the bracket is well outside the region where the
true solution lives, and is filtered out at the outer level by the
Armijo test.

\subsection{Convergence criteria}

A Newton iteration is declared converged when \emph{all} of the
following hold:
\begin{align}
  \lVert \mathsf{F}(\mathsf{x}_{k}) \rVert_{2} / \lVert \mathsf{x}_{k} \rVert_{2}
     &\;\le\; \texttt{nonlinearTolerance}, \\
  \lVert \mathsf{V}_{k} - \mathsf{V}_{k-1} \rVert_{2} /
     \lVert \mathsf{V}_{k} \rVert_{2}
     &\;\le\; \texttt{nonlinearVmTolerance}, \\
  \lVert \Iionsub{k} - \Iionsub{k-1} \rVert_{2} /
     \lVert \Iionsub{k} \rVert_{2}
     &\;\le\; \texttt{nonlinearIionTolerance}, \\
  \text{(opt.)}\quad
  \max_{i}\;
  \lVert \mathbf{s}_{k}^{(i)} - \mathbf{s}_{k-1}^{(i)} \rVert_{2} /
     \lVert \mathbf{s}_{k}^{(i)} \rVert_{2}
     &\;\le\; \texttt{nonlinearStatesTolerance},
\end{align}
and the iteration index satisfies $k + 1 \ge $
\code{minNonlinearIterations}. The last criterion is gated by the
boolean \code{nonlinearRequireStatesConvergence}. The coupled residual
$\lVert \mathsf{F} \rVert / \lVert \mathsf{x} \rVert$ has a noise floor
of order $\eta = $~\code{relTol} (RKF45) because $\mathsf{F}$ itself is
computed in finite precision through the ODE integrator. Tolerances
below this floor are unattainable in practice.

\subsection{Rollback on non-convergence}

If the Newton loop exits with the maximum iteration budget
\code{nonlinearIterations} reached and the convergence flag still
false, the boolean \code{nonlinearAcceptUnconverged} decides between:
\begin{itemize}[topsep=0pt]
  \item \code{false} (default): restore $\Vm, \Iion, \mathbf{s}$ to
        their values at the start of the time step, emit a
        \code{WarningInFunction}, and proceed; the time step
        effectively becomes a no-op.
  \item \code{true}: commit the unconverged iterate and proceed
        anyway.
\end{itemize}

\subsection{Alternative nonlinear methods}

For comparison and stress-testing, the solver also implements two
alternatives to JFNK selectable by \code{nonlinearMethod}:
\begin{description}[style=nextline]
  \item[\code{picard}] Fixed-point on the lagged Iion source: each
        outer iteration assembles
        $\Aimp\, \mathsf{x}^{(k+1)} = \mathsf{r}(\Iion(\mathsf{x}^{(k)}))$
        and solves a single KSP; the update is damped by
        \code{nonlinearRelaxation}. No matrix-free Jacobian and no
        line search are used. Robust at small $\Delta t$,
        linearly convergent.
  \item[\code{diagonalIion} (or \code{diagonal},
        \code{localDiagonal})] Picard with the diagonal of
        $\partial \Iion / \partial \Vm$ folded into the LHS
        before each linear solve. The diagonal derivative is
        evaluated by a one-sided finite difference of step
        \code{diagonalIionEpsilon}. The RHS is shifted
        accordingly so that the fixed point is preserved. Faster
        than pure Picard for stiff sources, no matrix-free
        Jacobian needed.
\end{description}

% ================================================================== %
\section{Linear solver: PETSc KSP}
% ================================================================== %

The wrapper \code{solveLinearKSP} drives any PETSc Krylov solver on
either the explicit operator $\Aimp$ (Picard / diagonal paths) or the
matrix-free Jacobian shell (JFNK path). The configuration is read from
\code{implicitHOCoeffs}:
\begin{center}
\begin{tabular}{@{}lll@{}}
\toprule
Key & Maps to & Default \\
\midrule
\code{linearSolver} & \code{KSPType} & \code{GMRES} \\
\code{preconditioner} & \code{PCType} & \code{ILU} \\
\code{tolerance} & relative tolerance & $10^{-10}$ \\
\code{maxIterations} & iteration cap & $2000$ \\
\code{jfnkMaxKrylovIterations} & GMRES restart cycle $m$ & $30$ \\
\code{jfnkMaxRestarts} & maximum outer restarts & $3$ \\
\bottomrule
\end{tabular}
\end{center}

\paragraph{Supported KSP types.} \code{GMRES}, \code{FGMRES},
\code{BiCGSTAB} (alias \code{BCGS}), \code{PREONLY} (alias
\code{SparseLU}, used together with a direct PC type).

\paragraph{Supported PC types.} \code{ILU}, \code{JACOBI},
\code{BJACOBI}, \code{LU} (alias \code{SparseLU}), \code{NONE},
\code{SOR}, \code{ASM}. Any preconditioner that PETSc supports is
also reachable via command-line overrides \code{-pc\_*} thanks to the
call to \code{KSPSetFromOptions}.

\paragraph{Orthogonalisation.} For \code{GMRES} and \code{FGMRES} the
wrapper explicitly selects \code{KSPGMRESModifiedGramSchmidtOrthogonalization},
which keeps the Krylov basis orthogonal to within $O(\varepsilon_{m} \kappa(\Aimp))$
instead of the classical $O(\varepsilon_{m} \kappa(\Aimp)^{2})$ \cite{saadbook},
relevant for the stiff Jacobians of cardiac upstrokes.

% ================================================================== %
\section{Ionic ODE integration}
% ================================================================== %

For each cell (or each cell quadrature point if \code{useHighOrder\_Iion}
is \code{true}) the TNNP ODE~\eqref{eq:tnnp} is advanced from $t^{n}$
to $t^{n+1}$ by the standard OpenFOAM adaptive integrator chosen via
the dictionary key \code{solver}. The default is \code{RKF45}, a 4-5
embedded Runge--Kutta--Fehlberg pair with PI step controller
\cite{hairer2}. The keys read from
\code{constant/timeIntegrationProperties} are:
\begin{center}
\begin{tabular}{@{}lll@{}}
\toprule
Key & Meaning & Default \\
\midrule
\code{solver} & Concrete ODE integrator & \code{RKF45} \\
\code{absTol} & Absolute tolerance per state & $10^{-9}$ \\
\code{relTol} & Relative tolerance per state & $10^{-6}$ \\
\code{maxSteps} & Cap on sub-steps per outer call & $10^{4}$ \\
\code{initialODEStep} & Initial sub-step size hint (ms) & $10^{-5}$ \\
\bottomrule
\end{tabular}
\end{center}

The relative tolerance \code{relTol} sets the noise floor in
$\mathsf{F}(\mathsf{x})$ and therefore the lowest attainable coupled
residual in the JFNK loop. Tightening \code{relTol} buys lower outer
tolerance at the cost of more RKF45 sub-steps per Newton iteration.
The cap \code{maxSteps} guards against runaway integrations triggered
by unphysical JFNK probes; combined with the Vm clamp of
section~5.4 it produces a clear \code{FatalError} when something is
truly wrong, rather than an indefinite hang.

% ================================================================== %
\section{TNNP numerical protection}
% ================================================================== %

The TNNP cell model contains several algebraic singularities
($i_{CaL}$ at $V = 0$ mV, divisions by gate variables that should be in
$[0,1]$, and stiff exponentials of unbounded $V$). The solver wraps
every external call into \code{TNNP::solveODE},
\code{TNNP::derivatives} and \code{TNNP::calculateCurrent} with a
\code{protectState} call that enforces physically reasonable bounds
\emph{before} the algebraic kernel is evaluated.

The master switch and per-guard knobs are read by the TNNP constructor
(\code{TNNP.C} lines 74--94) from the merged solver dictionary:

\begin{center}
\begin{tabular}{@{}lll@{}}
\toprule
Key & Type & Default \\
\midrule
\code{TNNPNumericalProtection} & Switch & \code{false} \\
\code{TNNPClampGates} & Switch & \code{true} \\
\code{TNNPClampVmForRates} & Switch & \code{true} \\
\code{TNNPLogNumericalProtection} & Switch & \code{true} \\
\code{TNNPConcentrationFloor} & scalar (mM) & $10^{-12}$ \\
\code{TNNPVMinForRates} & scalar (mV) & $-200$ \\
\code{TNNPVMaxForRates} & scalar (mV) & $+100$ \\
\code{TNNPVoltageZeroEps} & scalar (mV) & $10^{-8}$ \\
\code{TNNPMaxProtectionLogEntries} & label & $10^{6}$ \\
\bottomrule
\end{tabular}
\end{center}

When \code{TNNPNumericalProtection} is \code{true}, \code{protectState}
applies the following checks in order, recording every correction in
the in-memory \code{ProtectionSummaryEntry} table:

\begin{enumerate}[topsep=0pt,itemsep=2pt]
  \item \textbf{Non-finite $V$}. If \code{!std::isfinite(state[V])},
        $V$ is reset to the resting potential $-86.2$ mV. Reason tag
        \code{nonfinite\_V}.
  \item \textbf{Vm clamping for rates}
        (only if \code{TNNPClampVmForRates}). If $V <
        $~\code{TNNPVMinForRates} clamp to that floor;
        if $V > $~\code{TNNPVMaxForRates} clamp to that ceiling.
        Reason tags \code{Vm\_below\_minimum}, \code{Vm\_above\_maximum}.
  \item \textbf{Avoid the $i_{CaL}$ singularity at $V = 0$}. If
        $\lvert V \rvert < $~\code{TNNPVoltageZeroEps}, push
        $V$ to $\pm$\code{TNNPVoltageZeroEps} preserving the
        original sign. Reason tag
        \code{avoid\_i\_CaL\_V\_zero\_singularity}.
  \item \textbf{Concentration floor}. For each ionic concentration in
        $\{K_{i}, Na_{i}, Ca_{i}, Ca_{SR}\}$: non-finite entries are
        reset to \code{TNNPConcentrationFloor}; values below the floor
        are pushed to it. Reason tags
        \code{nonfinite\_concentration}, \code{concentration\_below\_floor}.
  \item \textbf{Gate clamping} (only if \code{TNNPClampGates}). For
        each gate in $\{Xr_{1}, Xr_{2}, Xs, m, h, j, d, f, fCa, s, r,
        g\}$: non-finite entries are reset to $0$; values $< 0$ are
        clamped to $0$; values $> 1$ are clamped to $1$. Reason tags
        \code{nonfinite\_gate}, \code{gate\_below\_zero},
        \code{gate\_above\_one}.
\end{enumerate}

\code{protectState} is invoked with a phase tag identifying its caller
(\code{calculateCurrent}, \code{solveODE\_start}, \code{solveODE\_end},
\code{derivatives}). All corrections at a given (phase, variable,
reason) triple are bucketed into a single
\code{ProtectionSummaryEntry} that tracks the total count, the first
correction time and integration point, and the min/max original values
and the maximum absolute delta. The summary is written on exit to
\code{TNNP\_numericalProtection\_summary.dat}.

% ================================================================== %
\section{Activation tracking and benchmarks}
% ================================================================== %

\subsection{Activation time per cell}

Once a cell crosses the activation threshold \code{activationThreshold}
($V_{\text{thr}}$) for the first time in a given time step, the
activation time is recorded by linear interpolation of $V$ across the
step,
\begin{equation}
  w \;=\; \mathrm{clip}\!\Bigl(
      \frac{V_{\text{thr}} - V^{n}}{V^{n+1} - V^{n}},
      \; 0, 1
    \Bigr),
   \qquad
   t_{\text{act}} \;=\; t^{n} + w\, \Delta t,
   \label{eq:tact}
\end{equation}
implemented in \code{updateActivationTimes}. A boolean flag
\code{calculateActivationTime[cellI]} is cleared after the first
crossing, so subsequent re-crossings (e.g.\ during recovery) are
ignored.

\subsection{Early termination on benchmark point}

When \code{stopAfterPointActivation} is \code{true}, the solver
identifies the cell nearest to \code{stopActivationPoint} (P8 in the
Niederer benchmark). As soon as that cell records a non-zero
$t_{\text{act}}$, the variable $\code{effectiveEndTime}$ is set to
$t_{\text{act}} + $~\code{stopDelayAfterActivation}, terminating the
simulation early once the wave has fully developed.

\subsection{Diagonal sampling}

Post-processing computes activation times at \code{nDiagonalSamples}
points uniformly distributed between \code{diagonalStart} and
\code{diagonalEnd} by inverse-distance-weighted interpolation
(\code{sampleIDW}):
\begin{equation}
  t_{\text{act}}(\mathbf{x}) \;=\;
   \frac{\sum_{j \in \mathcal{N}_{k}(\mathbf{x})}
         w_{j}\, t_{\text{act}, j}}
        {\sum_{j \in \mathcal{N}_{k}(\mathbf{x})} w_{j}},
   \qquad
   w_{j} \;=\; \frac{1}{\lVert \mathbf{x} - \mathbf{C}_{j} \rVert^{p}},
\end{equation}
where $\mathcal{N}_{k}$ is the set of \code{sampleNeighbours} nearest
cells and $p = $~\code{samplePower}. Default $k = 8$, $p = 2$.

% ================================================================== %
\section{Solver flow chart}
% ================================================================== %

The pseudo-code below summarises the per-time-step flow of the solver
and shows the \pyconst{run\_convergence.py constants}~/ \dictkey{dict
keys} that control each branch.

\begin{algorithm}[H]
\caption{High-level time-step loop of
         \texttt{highOrderElectroActivationFoamImplicitJfNK\_PETSc}.}
\begin{algorithmic}[1]
\State \texttt{PetscInitialize} (after \texttt{setRootCaseLists} so
       OpenFOAM owns MPI).
\State Read dictionaries.
\State Assemble PETSc \code{Mat}s: $\Mmass$ (\dictkey{massMatrix}),
       $\Kstif$ (\dictkey{useHighOrder\_Vm}),
       $\Lop = \Kstif/(\chi C_{m})$.
\State $\mathsf{V}^{0} \gets -84$~mV (initial condition).
\State $\mathsf{V}^{-1} \gets \mathsf{V}^{0}$,\ \,
       $\Delta t_{\text{prev}} \gets 0$,\ \,
       \texttt{hasPrev} $\gets$ false.
\While{ $t^{n} < $~\pyconst{END\_TIME} }
   \State $\Delta t \gets $~\pyconst{DT\_VALUES} (current sweep value).
   \State Build $\Aimp = \Mmass/\Delta t - \theta\,\Lop$,
          $\Bimp = \Mmass/\Delta t + (1-\theta)\,\Lop$
          via \code{MatDuplicate} + \code{MatScale} + \code{MatAXPY}
          (\dictkey{implicitScheme} $\to \theta$).
   \State Apply stimulus to $\Iext$ (\dictkey{stimulusProtocol}).
   \If{ \dictkey{jfnkInitGuessOrder} $\ge 1$ \textbf{and} \texttt{hasPrev} }
      \State $\mathsf{V}_{\text{guess}} \gets \mathsf{V}^{n}
             + (\Delta t/\Delta t_{\text{prev}})\,(\mathsf{V}^{n} - \mathsf{V}^{n-1})$,
             clamp to $[\,$\dictkey{jfnkInitGuessVmMin},
             \dictkey{jfnkInitGuessVmMax}$\,]$.
   \Else
      \State $\mathsf{V}_{\text{guess}} \gets \mathsf{V}^{n}$.
   \EndIf
   \State $\mathsf{V}^{n-1} \gets \mathsf{V}^{n}$,\ \,
          $\Delta t_{\text{prev}} \gets \Delta t$,\ \,
          \texttt{hasPrev} $\gets$ true.
   \If{ \dictkey{nonlinearMethod} = JFNK }
      \State $\mathsf{x} \gets \mathsf{V}_{\text{guess}}$.
      \For{ $k = 0, 1, \dots, $~\pyconst{IMPLICIT\_NONLINEAR\_ITERATIONS}$-1$ }
         \State $\mathsf{F}_{k} \gets$ \Call{computeResidual}{$\mathsf{x}$}
                (advances TNNP via \dictkey{ODE\_REL\_TOL},
                \dictkey{ODE\_ABS\_TOL}; clamp Vm by
                \dictkey{jfnkClampODEInput}).
         \If{ all convergence tests pass (\dictkey{nonlinearTolerance}, $\dots$) }
            \State \textbf{break}
         \EndIf
         \State Build a PETSc \code{MatShell} wrapping the FD
                Jacobian--vector product
                $\mathsf{J}\mathbf{v} \approx (\mathsf{F}(\mathsf{x}+\varepsilon\mathbf{v})-\mathsf{F}_{k})/\varepsilon$
                with $\varepsilon$ from \dictkey{jfnkEpsilon}.
         \State Solve $\mathsf{J}\,\boldsymbol{\delta} = -\mathsf{F}_{k}$
                with \code{KSPGMRES} + MGS (\dictkey{jfnkMaxKrylovIterations},
                \dictkey{jfnkLinearTolerance}, \dictkey{jfnkMaxRestarts}).
         \If{ \dictkey{jfnkLineSearch} }
            \State Armijo backtracking (\dictkey{jfnkArmijoC},
                   \dictkey{jfnkLineSearchMaxIter},
                   \dictkey{jfnkLineSearchAlphaMin}):
                   find $\alpha$ satisfying
                   $\lVert\mathsf{F}(\mathsf{x}+\alpha\boldsymbol{\delta})\rVert^{2}
                    \le (1-2c\alpha)\lVert\mathsf{F}_{k}\rVert^{2}$.
         \Else
            \State $\alpha \gets $ \dictkey{nonlinearRelaxation}.
         \EndIf
         \State $\mathsf{x} \gets \mathsf{x} + \alpha\,\boldsymbol{\delta}$.
      \EndFor
      \State $\mathsf{V}^{n+1} \gets \mathsf{x}$.
   \ElsIf{ \dictkey{nonlinearMethod} = picard or diagonalIion }
      \For{ $k = 0, 1, \dots, $~\pyconst{IMPLICIT\_NONLINEAR\_ITERATIONS}$-1$ }
         \State Solve $\Aimp\,\mathsf{V}_{*} = \mathsf{r}(\Iion(\mathsf{x}_{k}))$
                with \code{solveLinearKSP} (\dictkey{linearSolver},
                \dictkey{preconditioner}, \dictkey{tolerance}).
         \State $\mathsf{x}_{k+1} \gets
                (1{-}\omega)\,\mathsf{x}_{k} + \omega\,\mathsf{V}_{*}$,
                $\omega = $~\dictkey{nonlinearRelaxation}.
         \If{ convergence tests pass } \textbf{break} \EndIf
      \EndFor
      \State $\mathsf{V}^{n+1} \gets \mathsf{x}_{k+1}$.
   \EndIf
   \If{ \textbf{not} converged \textbf{and not} \dictkey{nonlinearAcceptUnconverged} }
      \State Warn; roll back: $\mathsf{V}^{n+1} \gets \mathsf{V}^{n}$,
             $\Iion^{n+1} \gets \Iion^{n}$, $\mathbf{s}^{n+1} \gets \mathbf{s}^{n}$.
   \EndIf
   \State Record activation times (\dictkey{activationThreshold}),
          P8 stop check (\dictkey{stopAfterPointActivation},
          \dictkey{stopDelayAfterActivation}).
   \State Write nonlinear-residual log if \dictkey{writeNonlinearResiduals}.
   \State $t^{n+1} \gets t^{n} + \Delta t$.
\EndWhile
\State Post-process: IDW sampling at P8 and diagonal
       (\dictkey{nDiagonalSamples}, \dictkey{sampleNeighbours},
        \dictkey{samplePower}).
\State \texttt{PetscFinalize} (after the inner \texttt{\{ \}} scope has
       destroyed all PETSc \code{Mat}s/\code{Vec}s).
\end{algorithmic}
\end{algorithm}

\noindent\textbf{Legend.} \pyconst{Blue} = Python constant set in
\code{run\_convergence.py}; \dictkey{green} = OpenFOAM dictionary key
read by the solver. Constants and keys are mapped one-to-one in
section~\ref{sec:driver}.

% ================================================================== %
\section{Driver script: \texttt{run\_convergence.py}}
\label{sec:driver}
% ================================================================== %

The Python driver \code{run\_convergence.py} writes the OpenFOAM
dictionaries and runs the solver in a parameter sweep. Every solver
knob is exposed as a top-level uppercase constant. The tables below
group the constants by topic; \emph{Key} is the dictionary entry
emitted by the script, and \emph{Section} points to the location in
this report where the underlying algorithm is described.

\subsection{Geometry and parameter sweep}
\begin{longtable}{@{}p{0.30\linewidth}p{0.28\linewidth}p{0.36\linewidth}@{}}
\toprule
Python constant & OpenFOAM key & Meaning \\
\midrule \endhead
\code{DIMENSIONS} & --- &
   List of dimensions to sweep (\code{2D}, \code{3D}). \\
\code{DX\_VALUES\_MM} & \code{blockMeshDict} cell counts &
   Uniform cell size in millimetres; mesh is generated by
   \code{blockMesh} from these. \\
\code{DT\_VALUES} & \code{controlDict.deltaT} &
   Outer time-step sizes to sweep (seconds). \\
\code{END\_TIME} & \code{controlDict.endTime} &
   Final simulation time (seconds). \\
\code{CFL} & \code{timeIntegrationProperties.CFL} &
   Informational reference Courant number; not enforced by
   the implicit solver. \\
\code{ACTIVE\_TIME\_SCHEMES} &
   \code{spatialIntegrationProperties.implicitScheme} &
   Subset of \{\code{backwardEuler}, \code{crankNicolson}\}. \\
\code{ACTIVE\_MASS\_MATRICES} &
   \code{spatialIntegrationProperties.massMatrix} &
   Subset of \{\code{lumped}, \code{consistent}\}. \\
\code{ACTIVE\_HO\_PAIRS\_Vm\_States} &
   \code{useHighOrder\_Vm} / \code{useHighOrder\_Iion}
   + \code{LRECoeffs\_*} & Tuples
   \texttt{(Vm-level, Iion-level)}; level is one of
   \code{NO}, \code{p1}, \code{p2}, \code{p3}. \\
\code{ACTIVE\_NONLINEAR\_METHODS} &
   \code{implicitHOCoeffs.nonlinearMethod} &
   Subset of \{\code{JFNK}, \code{picard}, \code{diagonalIion}, \dots\}. \\
\code{ACTIVE\_STATE\_ODE\_SOLVERS} &
   \code{timeIntegrationProperties.solver} &
   Subset of OpenFOAM ODE solver names. \\
\code{SIGMA\_X, SIGMA\_Y, SIGMA\_Z} & \code{cardiacProperties} &
   Diagonal components of the conductivity tensor $\Dten$. \\
\code{CHI} & \code{cardiacProperties.chi} & Surface-to-volume ratio. \\
\code{CM} & \code{cardiacProperties.Cm} & Membrane capacitance per area. \\
\bottomrule
\end{longtable}

\subsection{Nonlinear solver (JFNK / Picard)}
\begin{longtable}{@{}p{0.35\linewidth}p{0.30\linewidth}p{0.30\linewidth}@{}}
\toprule
Python constant & OpenFOAM key & Meaning \\
\midrule \endhead
\code{IMPLICIT\_NONLINEAR\_ITERATIONS} &
   \code{nonlinearIterations} &
   Hard cap on Newton/Picard outer iterations per time step. \\
\code{IMPLICIT\_MIN\_NONLINEAR\_ITERATIONS} &
   \code{minNonlinearIterations} &
   Minimum iterations before the convergence test is allowed
   to terminate the loop. \\
\code{NONLINEAR\_TOLERANCE} & \code{nonlinearTolerance} &
   Tolerance on the coupled residual
   $\lVert \mathsf{F} \rVert / \lVert \mathsf{x} \rVert$.
   Cannot be below the ODE relTol floor. \\
\code{NONLINEAR\_VM\_TOLERANCE} & \code{nonlinearVmTolerance} &
   Relative L$^{2}$ change in $\Vm$ between successive iterates. \\
\code{NONLINEAR\_IION\_TOLERANCE} & \code{nonlinearIionTolerance} &
   Same metric on the $\Iion$ field. \\
\code{NONLINEAR\_STATES\_TOLERANCE} &
   \code{nonlinearStatesTolerance} &
   Maximum over state index of relative L$^{2}$ change in
   each ionic state. \\
\code{NONLINEAR\_REQUIRE\_STATES\_CONVERGENCE} &
   \code{nonlinearRequireStatesConvergence} &
   If \code{true}, the states test is required in addition
   to the others. \\
\code{NONLINEAR\_RELAXATION} & \code{nonlinearRelaxation} &
   Initial Newton step length $\alpha_{0}$ (also the
   Picard relaxation factor in non-JFNK branches). \\
\bottomrule
\end{longtable}

\subsection{JFNK matrix-free options}
\begin{longtable}{@{}p{0.35\linewidth}p{0.30\linewidth}p{0.30\linewidth}@{}}
\toprule
Python constant & OpenFOAM key & Meaning \\
\midrule \endhead
\code{JFNK\_MAX\_KRYLOV\_ITERATIONS} &
   \code{jfnkMaxKrylovIterations} &
   GMRES restart cycle $m$. \\
\code{JFNK\_LINEAR\_TOLERANCE} & \code{jfnkLinearTolerance} &
   Relative tolerance of the inner KSP. \\
\code{JFNK\_EPSILON} & \code{jfnkEpsilon} &
   Forward-difference perturbation $\varepsilon_{0}$ in
   \eqref{eq:fdjac}. Set near $\sqrt{\eta}$, where $\eta$ is the
   ODE \code{relTol}. \\
\code{DIAGONAL\_IION\_EPSILON} & \code{diagonalIionEpsilon} &
   Step size for the diagonal $\partial \Iion / \partial \Vm$
   in the \code{diagonalIion} branch. \\
\bottomrule
\end{longtable}

\subsection{Linear solver}
\begin{longtable}{@{}p{0.30\linewidth}p{0.30\linewidth}p{0.35\linewidth}@{}}
\toprule
Python constant & OpenFOAM key & Meaning \\
\midrule \endhead
\code{IMPLICIT\_LINEAR\_SOLVER} & \code{linearSolver} &
   PETSc KSP type used by the explicit-operator solve
   (Picard / diagonal branches). \\
\code{IMPLICIT\_TOLERANCE} & \code{tolerance} &
   Relative residual tolerance of the explicit-operator KSP. \\
\code{IMPLICIT\_MAX\_ITERATIONS} & \code{maxIterations} &
   Iteration cap on the explicit-operator KSP. \\
\bottomrule
\end{longtable}

\subsection{Inner ODE integrator (RKF45)}
\begin{longtable}{@{}p{0.30\linewidth}p{0.30\linewidth}p{0.35\linewidth}@{}}
\toprule
Python constant & OpenFOAM key & Meaning \\
\midrule \endhead
\code{ODE\_ABS\_TOL} & \code{timeIntegrationProperties.absTol} &
   Absolute error tolerance per state. \\
\code{ODE\_REL\_TOL} & \code{timeIntegrationProperties.relTol} &
   Relative error tolerance per state. Drives the noise floor
   of the matrix-free Jacobian. \\
\code{ODE\_INITIAL\_STEP} & \code{initialODEStep} &
   Initial sub-step size hint. \\
\code{ODE\_MAX\_STEPS} & \code{maxSteps} &
   Cap on RKF45 sub-steps per outer call. Above this
   the solver raises a fatal error. \\
\bottomrule
\end{longtable}

\subsection{LRE high-order parameters per p-level}

\code{P\_LEVELS\_BY\_DIM[dim][p]} is a dictionary giving, for each
combination of geometric dimension and polynomial order:
\begin{center}
\begin{tabular}{@{}lll@{}}
\toprule
Sub-key & Meaning & Typical value \\
\midrule
\code{N} & Polynomial order ($1$, $2$ or $3$) & matches the \code{p} tag \\
\code{Nn} & Number of stencil neighbours required & 13--60 depending on (dim, p) \\
\code{k} & Stencil search radius (in cell diameters) & 6 \\
\code{maxStencilSize} & Hard cap on stencil size & 33--80 \\
\bottomrule
\end{tabular}
\end{center}
The script writes these into a nested
\code{highOrderCoeffs\{LRECoeffs\_Vm\{\dots\}; LRECoeffs\_Iion\{\dots\}\}}
block in \code{spatialIntegrationProperties} whenever HO is enabled
for either Vm or $\Iion$.

\subsection{Post-processing and run-control}
\begin{longtable}{@{}p{0.35\linewidth}p{0.30\linewidth}p{0.30\linewidth}@{}}
\toprule
Python constant & OpenFOAM key & Meaning \\
\midrule \endhead
\code{ACTIVATION\_THRESHOLD} & \code{activationThreshold} &
   $V_{\text{thr}}$ for activation-time crossing (V). \\
\code{N\_DIAGONAL\_SAMPLES} & \code{nDiagonalSamples} &
   Number of sample points along the diagonal. \\
\code{SAMPLE\_NEIGHBOURS} & \code{sampleNeighbours} &
   Number of nearest cells used by IDW. \\
\code{SAMPLE\_POWER} & \code{samplePower} &
   Inverse-distance power $p$. \\
\code{STOP\_AFTER\_POINT\_ACTIVATION} &
   \code{stopAfterPointActivation} &
   Terminate after the P8 cell activates. \\
\code{STOP\_DELAY\_AFTER\_ACTIVATION} &
   \code{stopDelayAfterActivation} &
   Wall-clock-independent delay after P8 activation before
   stopping (seconds). \\
\code{WRITE\_NONLINEAR\_RESIDUALS} &
   \code{writeNonlinearResiduals} &
   Emit per-iteration residual log to
   \code{postProcessing/.../nonlinearResiduals.dat}. \\
\code{SAVE\_LOGS\_ON\_SUCCESS} & --- &
   Python-only flag; if \code{false}, deletes solver and
   blockMesh log files after a successful run. \\
\bottomrule
\end{longtable}

% ================================================================== %
\appendix
\section{Dimension table}
% ================================================================== %

OpenFOAM dimension vectors (mass, length, time, temperature, amount,
current, luminous intensity) of the fields declared in
\code{createFields.H}:

\begin{center}
\begin{tabular}{@{}lll@{}}
\toprule
Field & OpenFOAM \code{dimensionSet} & SI \\
\midrule
\code{Vm} & $[1, 2, -3, 0, 0, -1, 0]$ & volt \\
\code{Iion} & \code{dimVoltage/dimTime} & \si{\volt\per\second} \\
\code{externalStimulusCurrent} & \code{dimCurrent/dimVolume} & \si{\ampere\per\meter\cubed} \\
\code{conductivity} & \code{dimCurrent\^{}2 dimTime\^{}3 / (dimMass dimVolume)} & \si{\siemens\per\meter} \\
\code{chi} & \code{dimless/dimLength} & \si{\per\meter} \\
\code{Cm} & \code{dimCurrent dimTime / (dimVoltage dimArea)} & \si{\farad\per\meter\squared} \\
\bottomrule
\end{tabular}
\end{center}

% ================================================================== %
\section{File map}
% ================================================================== %

\begin{center}
\begin{tabular}{@{}p{0.55\linewidth}p{0.40\linewidth}@{}}
\toprule
Path (relative to \code{cardiacFoam-pc/}) & Role \\
\midrule
\code{applications/solvers/.../JfNK\_PETSc.C} & Main translation unit; time loop, JFNK loop, KSP wrapper. \\
\code{applications/solvers/.../createFields.H} & Field instantiation and dictionary reading. \\
\code{src/ionicModels/TNNP/TNNP.H, .C} & TNNP cell model and protection guards. \\
\code{src/ionicModels/TNNP/TNNP\_2004.H, *Names.H} & Generated CellML kernel and state-name table. \\
\code{src/ionicModels/ionicModel/ionicModel.{H,C}} & Base class with run-time selection of the ODE solver. \\
\code{tutorials\_electro/.../run\_convergence.py} & Python driver of the parameter sweep. \\
\code{constant/cardiacProperties} & $\chi$, $C_{m}$, conductivity tensor, ionic model selection. \\
\code{constant/timeIntegrationProperties} & ODE solver, tolerances, sub-step cap. \\
\code{constant/spatialIntegrationProperties} & Time scheme, HO flags, JFNK / KSP / line-search options. \\
\code{constant/stimulusProtocol} & Stimulus boxes, intensity, duration, start times. \\
\bottomrule
\end{tabular}
\end{center}

% ================================================================== %
\begin{thebibliography}{9}
% ================================================================== %

\bibitem{ttp2004} K.~H.~W.~J. ten Tusscher, D. Noble, P.~J. Noble, A.~V.
    Panfilov, \emph{A model for human ventricular tissue},
    American Journal of Physiology -- Heart and Circulatory Physiology
    \textbf{286} (2004) H1573--H1589.

\bibitem{niederer2012} S.~A. Niederer et al.,
    \emph{Verification of cardiac tissue electrophysiology simulators
    using an N-version benchmark}, Progress in Biophysics and Molecular
    Biology \textbf{104} (2012) 22--48.

\bibitem{knoll2004} D.~A. Knoll, D.~E. Keyes,
    \emph{Jacobian-Free Newton-Krylov methods: a survey of approaches
    and applications}, Journal of Computational Physics \textbf{193}
    (2004) 357--397.

\bibitem{noceWright} J. Nocedal, S.~J. Wright,
    \emph{Numerical Optimization}, 2nd edition, Springer, 2006
    (Chapter 3, line-search methods).

\bibitem{saadbook} Y. Saad,
    \emph{Iterative Methods for Sparse Linear Systems}, 2nd edition,
    SIAM, 2003.

\bibitem{hairer2} E. Hairer, G. Wanner,
    \emph{Solving Ordinary Differential Equations II: Stiff and
    Differential-Algebraic Problems}, Springer, 1996.

\end{thebibliography}

\end{document}
